\section{Method}
This chapter describes the methods used to optimize the Majorana modes in the junction setup and to analyze their properties. We first discuss the parameter optimization process, including the model parameterization, loss function design based on Majorana criteria, and the optimization strategy. Next, we detail the post-optimization analysis of the setup, covering Hamiltonian reconstruction, parity classification, and various diagnostics. Finally, we describe the braiding protocol implemented in an effective Majorana model and how we validate the results using gate-validation checks.
\subsection{Parameter Optimization}
A key part of this thesis is the numerical search for parameter regimes in which an interacting $n$-dot PMM system exhibits Majorana-like low-energy states. For the non-interacting two-dot problem, the sweet-spot conditions are known analytically, but once Coulomb interactions, longer chains, and constraints on higher excited states are included, the problem quickly becomes too complicated to solve analytically. We therefore treat the search for suitable parameters as an optimization problem, where the Hamiltonian is evaluated repeatedly and compared against a set of Majorana criteria. \\ \\
The purpose of the optimization is not to claim a unique exact sweet spot for a given device. Instead, it is to identify parameter regions that repeatedly produce near-degenerate, well-localized, and spectrally protected Majorana-like states, and to study how these regions change with the number of dots and with the interaction strength. In the calculations presented here, the Coulomb repulsion $U$ is fixed to selected values, while the remaining free parameters are optimized within that interaction regime.
\subsubsection{Model parameterization}
The model parameterization is simply based on the n-dot PMM Hamiltonian:
\begin{equation}
H = \sum_{i=1}^{n} \epsilon_i n_i + \sum_{i=1}^{n-1} \left[ -t_i \left(c_i^{\dagger} c_{i+1} + c_{i+1}^{\dagger} c_i\right) + \Delta_i \left(c_i c_{i+1} + c_{i+1}^{\dagger} c_i^{\dagger}\right) + U n_i n_{i+1} \right]
\end{equation}

For an $n$-dot system, this Hamiltonian contains $n$ on-site energies $\epsilon_i$, $n-1$ hopping amplitudes $t_i$, $n-1$ pairing amplitudes $\Delta_i$, and nearest-neighbor Coulomb terms. The optimization is carried out on a reduced parameter vector containing only the free parameters. Depending on the chosen ansatz, a parameter can either be treated as homogeneous, inhomogeneous, or fixed to a prescribed value. This makes it possible to compare different levels of model complexity while keeping the numerical procedure the same. \\ \\
In the calculations used throughout this thesis, the starting point is chosen close to the known non-interacting sweet spot, with $t_i \approx \Delta_i$ and on-site energies near the corresponding sweet-spot values. The interaction strength $U$ is then fixed, and the remaining parameters are allowed to adjust so that the optimizer searches for the best Majorana-like configuration within that interaction regime.
\subsubsection{Loss function and Majorana criteria}
The loss function is designed to optimize the parameters of the n-dot PMM Hamiltonian such that the resulting system exhibits properties consistent with Majorana modes. The loss function is based on several criteria that are commonly associated with stable and protected Majorana modes:\\ \\
\textbf{Energy degeneracy}: Energy degeneracy is highly weighted in the loss function, especially the degeneracy of the ground state. The weight of the energy degeneracy is based on which energy level we are looking at. For example, the degeneracy of the ground state is more important than the degeneracy of the first excited state, and so on. This is because the ground-state degeneracy is a key signature of Majorana modes, while degeneracies in higher energy levels may not be as robust or relevant for topological protection. The array of weights linearly decreases from some weight $\omega_\text{max}$ for the ground-state degeneracy to a smaller weight $\omega_\text{min}$ for the degeneracy of the highest energy level considered in the loss function. This way, we can prioritize the optimization towards achieving a degenerate ground state while still considering the overall energy spectrum.\\
\textbf{Degeneracy Gaps}: The loss function also includes terms that penalize small energy gaps between energy levels. This is not as highly weighted as the energy degeneracy, but it is still important to ensure that there are sufficiently large energy gaps to protect the Majorana modes from thermal excitations and other perturbations. The weighting of the energy gaps is based on the same idea as the energy degeneracy, where the weight linearly decreases from some weight $\omega_\text{max}$ for the gap between the ground state and the first excited state, to a smaller weight $\omega_\text{min}$ for the gap between the highest energy levels considered in the loss function. We want clearly separated energy levels while performing a braiding protocol. Better energy gaps help to ensure that the system remains in the desired energy subspace during the braiding process, which is crucial for maintaining the topological protection of the Majorana modes.\\
\textbf{Charge difference}: The function also includes a term that penalizes charge differences between the degenerate states. This is because, when Majorana modes are present, the degenerate states should have the same charge. A significant charge difference between the degenerate states would indicate that they are not true Majorana modes. \\
\textbf{Majorana Polarization}: On the edge dots of the PMM system, the Majorana polarization should equal $\pm 1$, while the inner dots should always be zero. The loss function calculates the Majorana polarization for each dot and includes terms that penalize deviations from these ideal values. This criterion is important because it reflects the spatial localization of the Majorana modes, which is a key feature for their stability and potential use in quantum computing applications. By optimizing the parameters to achieve the desired Majorana polarization profile, we can enhance the likelihood of realizing robust Majorana modes in the system.\\ \\
If all of the above criteria are satisfied, the system is likely to host Majorana modes with desirable properties. In practice, the different penalty terms are combined into a single weighted loss and normalized by the mean energy scale of the spectrum, so that different parameter sets can be compared on a more equal footing.

\subsubsection{Optimization strategy}\label{sec:optimization_strategy}
The numerical search is carried out in two stages. First, a local gradient-based optimization is performed using the Adam algorithm. The optimizer acts on the reduced parameter vector, while after each update the vector is mapped back to the physical parameter set. In this step, fixed parameters are reinserted and constraints such as positive couplings are enforced. The local optimization therefore explores only physically allowed Hamiltonians, while still benefiting from automatic differentiation of the loss function. \\ \\
Because the loss landscape contains many local minima, a single local run is usually not sufficient. We therefore supplement the local optimizer with repeated restarts and a basin-hopping procedure. Starting from a trial parameter set, several locally optimized runs are performed from randomly perturbed initial conditions in its neighborhood, and the lowest-loss result is retained as the representative of that basin. This improves the robustness of the search and reduces the risk that the final result depends too strongly on a single initialization. \\ \\
To move between basins, a random step is added to the current parameter vector before the local refinement is repeated. The new candidate is accepted whenever it improves the loss, and can also be accepted with a Metropolis probability when it is slightly worse. In this way, the search can occasionally leave a shallow local minimum and continue exploring other regions of parameter space. For each choice of $n$ and fixed $U$, the parameter set with the lowest final loss is kept for the subsequent spectral analysis and braiding calculations.

\subsection{Post-Optimization Analysis of the Setup}
\subsubsection{Hamiltonian reconstruction and parity classification}
Once an optimized parameter set has been found, the raw optimization vector is converted back into physical parameters and the full many-body Hamiltonian is reconstructed. This step is important because the optimizer itself only works with a reduced parameter vector, whereas the physical interpretation requires the complete set of hopping, pairing, on-site-energy, and interaction parameters. In practice, fixed parameters are reinserted, homogeneous parameters are expanded where needed, and the Hamiltonian is rebuilt in the same basis used during optimization. For selected configurations, the parameter values can also be exported and reloaded in a separate analysis setup, where the Hamiltonian is reconstructed numerically and, when useful, symbolically. \\ \\
The reconstructed Hamiltonian is then diagonalized to obtain its full spectrum and eigenstates. To identify the low-energy structure relevant for Majorana physics, the eigenstates are classified according to the total fermion parity operator,
\[
P_{\mathrm{tot}} = \prod_{i=1}^{n} (1 - 2n_i).
\]
This separates the spectrum into even- and odd-parity sectors, which makes it possible to compare candidate degenerate partners level by level. The parity-resolved spectrum is used to inspect the ground-state splitting, the excitation gap above the low-energy manifold, and the overall pairing of even and odd states. This classification step also defines the low-energy basis that is used later when extracting effective Majorana operators and constructing the braiding protocol.
\subsubsection{Spectral and local diagnostics}
The optimization condenses several physical criteria into a single scalar loss, but that loss alone does not explain why a parameter set is good. For this reason, each promising configuration is analyzed using a set of complementary diagnostics. At the spectral level, we inspect the parity-resolved energy levels, the total even-odd splitting, and the charge difference between candidate degenerate partners. In addition, the optimized parameters are visualized together with the spectrum to relate spectral features directly to the physical setup. \\ \\
At the operator level, we compute observables that probe how Majorana-like the low-energy states actually are. These include the Majorana polarization profile across the chain, the local electron occupancy in even and odd states, the matrix elements of local Majorana operators between parity partners, and the hopping and pairing expectation values on each bond. In a separate analysis pass, we also evaluate transition matrices
\[
T_{ss'}^{(O)} = \langle s \vert O \vert s' \rangle
\]
for relevant operators such as local charge operators, local Majorana operators, and nearest-neighbor hopping and pairing terms. These diagnostics help distinguish configurations that merely achieve a low numerical loss from configurations that exhibit a clear low-energy Majorana interpretation. Only after this analysis step do we select configurations for the braiding calculations described in the following sections.

\subsection{Braiding in an Effective Majorana Model}
\subsubsection{Pulse protocol}
Before applying the braiding procedure to the optimized microscopic setup, we first test it in an idealized effective model containing four Majorana operators $γ_0$, $γ_1$, $γ_2$, and $γ_3$. The purpose of this simplified model is to isolate the braiding protocol itself from the complications of the full interacting Hamiltonian. In this setting, the time-dependent Hamiltonian is written as
\[
H(t) = \Delta_1(t)\, i γ_0 γ_1 + \Delta_2(t)\, i γ_0 γ_2 + \Delta_3(t)\, i γ_0 γ_3,
\]
where the couplings $\Delta_1(t)$, $\Delta_2(t)$, and $\Delta_3(t)$ are turned on and off sequentially through smooth pulses. The pulses are chosen as broadened step-like functions with controllable width and steepness, so that the exchange path remains continuous in parameter space.

The braiding cycle is implemented by keeping one coupling active at the beginning and end of the protocol, while the two remaining couplings are activated at intermediate times. In this way, the dominant Majorana pairing is moved through the network in a prescribed order. At each time step, the instantaneous Hamiltonian is diagonalized and its spectrum is monitored to verify that the ground-state manifold remains approximately degenerate and separated from higher excited states throughout the cycle.

\subsubsection{Kato Berry-phase calculation}
To determine which gate is implemented by the braiding cycle, we compute the non-Abelian Berry phase of the instantaneous low-energy subspace using the Kato formulation of adiabatic transport. Starting from the degenerate ground-state manifold at $t=0$, we construct the projector
\[
P(t) = \sum_{a \in \mathrm{GS}} \ket{\psi_a(t)}\bra{\psi_a(t)},
\]
and numerically propagate the associated Kato evolution operator along the discretized path in time. In practice, this is done by diagonalizing the Hamiltonian at each time step, constructing the projector onto the low-energy manifold, and updating the transport operator from the change in the projector between successive steps.

This procedure yields a unitary operator acting on the degenerate subspace at the end of the loop. The resulting unitary can then be compared with the expected Majorana exchange operator. In the effective model, this provides a clean reference for what the ideal braid should do before the same protocol is applied to a more realistic projected Hamiltonian.

\subsection{Braiding with Optimized Majorana Modes}
\subsubsection{Extraction of Majorana operators}
After identifying optimized parameter sets, we construct effective Majorana operators directly from the low-energy even- and odd-parity eigenstates of each subsystem. For each subsystem, the even and odd states are paired according to their energy proximity, and operator combinations of the form
\[
\gamma_1 \sim \ket{o}\bra{e} + \ket{e}\bra{o}, \qquad
\gamma_2 \sim i(\ket{o}\bra{e} - \ket{e}\bra{o})
\]
are formed. By summing over a chosen number of low-energy parity pairs, we obtain effective Majorana operators that act within the relevant low-energy sector.

These subsystem operators are then embedded into the full multi-arm Hilbert space using Jordan-Wigner strings, so that the resulting Majorana operators satisfy the correct fermionic anticommutation structure across the entire device. In this way, the idealized Majorana degrees of freedom used in the toy model are replaced by operators extracted from the optimized microscopic Hamiltonian.

\subsubsection{Projected microscopic junction Hamiltonian}
To move beyond the idealized braid Hamiltonian, we construct junction operators that couple neighboring subsystems through microscopic hopping and pairing terms. These operators are first defined in the full Hilbert space and are then projected onto the low-energy manifold spanned by the relevant optimized states. The projected operators,
\[
T_{AB} = P^\dagger H^{\mathrm{junc}}_{AB} P,
\qquad
T_{AC} = P^\dagger H^{\mathrm{junc}}_{AC} P,
\]
replace the hand-chosen Majorana bilinears of the toy model.

The time-dependent braiding Hamiltonian is then built by driving these projected junction operators with the same pulse sequence used in the effective model. This keeps the geometry of the braid fixed while allowing us to test whether the microscopic optimized setup reproduces the desired exchange dynamics in its low-energy sector.

\subsubsection{Gate-validation checks}
The final step is to test whether the projected microscopic braid behaves like the intended Majorana exchange. Several consistency checks are performed along the braiding path. First, we verify basic structural properties such as Hermiticity, parity conservation, ground-state splitting, and the excitation gap above the low-energy manifold. Next, we test the algebra of the extracted Majorana operators by checking Hermiticity, normalization, and mutual anticommutation.

After the Kato evolution has been computed, the resulting unitary is compared with the expected exchange action on the Majorana operators. In particular, we test whether a single exchange maps $γ_2 \to -γ_3$ and $γ_3 \to γ_2$, whether a double exchange produces the expected sign change, and whether repeated exchanges return the system to its initial operator structure. We also project the braid unitary into parity sectors and compare the resulting blocks with the target gate. Together, these checks determine whether the optimized setup supports a braid that is not only spectrally well-defined, but also implements the intended topological operation.
